\documentclass{article}
\usepackage[T1]{fontenc}
\usepackage{lmodern}
\usepackage{graphicx}
\usepackage{amsmath,amssymb}
\usepackage[a4paper,margin=2cm]{geometry}
\usepackage[absolute,overlay]{textpos}
\setlength{\TPHorizModule}{1mm}
\setlength{\TPVertModule}{1mm}
\usepackage[indonesian]{babel}
\usepackage{hyperref}
\setlength{\emergencystretch}{2em}
% unit: Halaman 18

\begin{document}

% === Halaman 18 (python/native) ===

18

• Jika panjang plainteks tidak habis dibagi dengan ukuran blok, maka blok 

terakhir berukuran lebih pendek daripada blok-blok lainnya. 

• Untuk itu, kita tambahkan bit-bit padding untuk menutupi kekurangan 

bit blok.

• Bit-bit padding misalnya bit 0 semua, atau bit 1 semua, atau bit 0 dan 

bit 1 berselang-seling.

• Contoh: Ukuran blok 8 bit: 10001101  00101001  10110000


Blok terakhir hanya 5 bit, ditambah 3 bit padding (warna biru)

\end{document}
